\begin{textsample}{POS Dim 4 – global\_north\_2025\_01 – Score 30.00 – global\_north\_2025\_01/120022352.txt}  \label{ex:f4_pos_048}
After 16 months, the current bombardment of Gaza and warfare between Hamas and Israel could soon be over, with Palestinian sources telling Andalou Agency a \textbf{ceasefire} \textbf{deal} could be signed by Friday.
The \textbf{draft} has been presented at \textbf{talks} in Doha involving Israel, Hamas, the US ( including representatives for both \textbf{President} Joe Biden and \textbf{President-elect} Donald Trump ), Qatar and Egypt.
By our count, this will have been at least the eighth \textbf{ceasefire} \textbf{proposal} since October 2023.
Here is a timeline of the previous \textbf{proposals}, and some of the reasons why they failed.
Nov 2023 : \textbf{Proposed} \textbf{extension} of \textbf{temporary} \textbf{ceasefire} The initial Qatar-mediated \textbf{temporary} \textbf{ceasefire} \textbf{agreement} secured a four-day pause in fighting, from November 24-28, to allow for the \textbf{release} of 50 Israeli hostages, 120 Palestinian political \textbf{prisoners}, and for humanitarian aid to enter Gaza.
A two-day \textbf{extension} was agreed to in \textbf{exchange} for the \textbf{release} of a further 20 Israelis and 60 Palestinians.
After a final one-day \textbf{extension} securing the \textbf{release} of eight Israeli hostages @ @ @ @ @ @ @ @ @ @ December 1, after \textbf{negotiations} to further \textbf{extend} failed due to violations by both Israel and Hamas.
Dec 2023 : US vetoes UN \textbf{ceasefire} resolution The United Arab Emirates \textbf{proposed} a resolution to the UN Security Council, calling for an immediate humanitarian \textbf{ceasefire}, unconditional \textbf{release} of all hostages, and the \textbf{resumption} of humanitarian access to Gaza.
The resolution stated both"Palestinian and Israeli civilian populations must be protected in accordance with international humanitarian law".
At a December 8 vote the resolution was supported by 13 council members but vetoed by the US ( using its power as one of five \textbf{permanent} members of the Security Council ).
The UK abstained from voting.
Both cited the failure to condemn Hamas as their reason for not supporting the resolution. \textbf{Feb} 2024 : Further vetoes at the UN On February 20, the US vetoed a \textbf{ceasefire} resolution presented by Algeria.
It contained similar \textbf{terms} to the previous UAE-drafted resolution, with additional statements on preventing further escalation in the Middle East.
The US claimed @ @ @ @ @ @ @ @ @ @ own \textbf{ceasefire} \textbf{negotiations} between Israel and Hamas.
It instead put forward an alternative resolution, which made a \textbf{ceasefire} conditional on the \textbf{release} of hostages.
This was vetoed by \textbf{permanent} Security Council members Russia and China, who argued any \textbf{ceasefire} should be unconditional.
Mar 2024 : Israel ignores successful UN resolution The fourth attempt by the Security Council to pass a \textbf{ceasefire} resolution succeeded, with the US abstaining from the vote.
This \textbf{proposal} was put forward by the non-permanent members of the council ( known as the E-10 or"elected 10") and called for a \textbf{ceasefire} for the duration of Ramadan as the first step towards a"lasting sustainable \textbf{ceasefire}".
Israel ' s then foreign minister Israel Katz immediately confirmed that Israel would not abide by the \textbf{ceasefire}, insisting his country would"destroy Hamas and continue to fight until the last of the hostages \textbf{returns} home".
May 2024 : Hamas \textbf{accepts} Egypt-Qatar \textbf{proposal}, Israel invades Rafah Stage one included a \textbf{temporary} \textbf{ceasefire} and retreat of Israeli forces, to @ @ @ @ @ @ @ @ @ @ of displaced Palestinians, \textbf{release} of Israeli hostages and Palestinian political \textbf{prisoners}, and the beginning of \textbf{reconstruction} in Gaza.
Stage two included the \textbf{permanent} end of military operations, Israel ' s complete \textbf{withdrawal} from Gaza, and the further \textbf{exchange} of \textbf{captive} Israelis and Palestinians.
Stage three included the final \textbf{exchange} of remaining \textbf{captive} Israelis and Palestinians, the beginning of a long-term \textbf{reconstruction} \textbf{plan} for Gaza, and the end of the siege on the Gaza Strip.
A few weeks later, the US put forward a very similar \textbf{proposal} -- with the key difference being it removed the reference to ending the siege on the Gaza Strip.
Hamas was reportedly also receptive to this \textbf{deal}.
In June 2024, Netanyahu \textbf{rejected} the \textbf{deal}, stating he was only interested in a"partial"\textbf{ceasefire} and was"committed to continuing the war after a pause."Aug 2024 : Hamas \textbf{rejects} ' bridging \textbf{proposal} ' In late August, the US, Egypt and Qatar revisit the \textbf{plan} and put forward a"bridging @ @ @ @ @ @ @ @ @ @ \textbf{ceasefire}"period, with a focus on the \textbf{exchange} of vulnerable Israeli and Palestinian \textbf{captives}, that could be \textbf{extended} indefinitely while all parties negotiated the details of the second \textbf{phase} of the \textbf{plan}.
Israel \textbf{accepted} this \textbf{proposal} but Hamas \textbf{rejected} it, wanting to \textbf{return} to the previous \textbf{three-phase} \textbf{plan}.
The partial \textbf{withdrawal} of Israeli forces from key areas in stage one, to allow safe passage of aid and displaced civilians to \textbf{return} to northern Gaza.
Hamas is to provide a list of Israeli \textbf{prisoners} currently in custody, allowing for an \textbf{exchange} of Palestinian \textbf{prisoners} for some \textbf{captive} Israeli soldiers.
An \textbf{agreement} on a 1,000-meter \textbf{buffer} zone. \textbf{Negotiations} to continue throughout \textbf{phase} one, to confirm details for \textbf{phases} two and three.
Have something to say about this article?
Write to us at letters@crikey.com.au.
Please include your full name to be considered for publication in Crikey'sYour Say.
We reserve the right to edit for length and clarity.
Share About the Author Crystal Andrews is Crikey ' s readers ' @ @ @ @ @ @ @ @ @ @, a digital media outlet producing feminist commentary on news, social issues, digital and pop culture for young Australian women.
Topics Crikey encourages robust conversations on our website.
However, we ' re a small team, so sometimes we have to reluctantly turn comments off due to legal risk.
You can read our moderation guidelines here.
If you ' d like to share your thoughts on this story, you can email letters@crikey.com.au.
Please include your full name to be considered for publication in our Your Say section.
Crikey is an independent Australian source for news, investigations, analysis and opinion focusing on politics, media, economics, health, international affairs, the climate, business, society and culture.
We are guided by a deceptively simple, old idea : tell the truth and shame the devil.

% matched lemmas: accept, agreement, buffer, captive, ceasefire, deal, draft, exchange, extend, extension, feb, negotiation, permanent, phase, plan, president, president-elect, prisoner, proposal, propose, reconstruction, reject, release, resumption, return, talk, temporary, term, three-phase, withdrawal
\end{textsample}
